\section{A Weil character whose perfect representative is obstructed}
\label{sec:equivariant-eightfolds}
This section separates the construction of a Weil class in $K$-theory from the deformation of a perfect complex representing it.

Perry's Theorem~1.2 \cite{Perry}, an unrefereed preprint, extends the semiregularity method used by Markman. In a family of abelian varieties it permits a perfect complex $E$ with $\operatorname{Ext}^{<0}(E,E)=0$, a finite group $G\subset A\times A^\vee$ acting by translations and degree-zero line bundles, and weak $G$-semiregularity. If the prescribed twisted Chern character remains Hodge, the theorem deforms $E$ as a twisted perfect complex and propagates algebraicity. An isolated CM automorphism is not automatically one of the permitted group actions. More fundamentally, a virtual class with the right Chern character need not have a semiregular representative.

We give an explicit calculation at the split Gaussian eightfold
\[
 A_0=X\times Y,
 \qquad X=Y=E_i^4,
 \qquad a(x,y)=(ax,\bar a y),
\]
with product principal polarization. Its Hermitian lattice is $I_{4,4}$. Let $\Delta\subset X\times Y$ be the diagonal. If $a=2+i$, write $S=a^*+\bar a^*$ for the sum of pullback operations on coherent $K$-theory, and put
\[
 P(s)=(s+1170)(s+350)(s-750)(s-1250).
\]

\begin{proposition}\label{prop:graph-obstruction}
The integral class $v=P(S)[\mathcal O_\Delta]$ has a perfect representative $E_v$ satisfying
\[
 \operatorname{ch}(E_v)=-339429556224\,e_W[\Delta]\ne0,
 \qquad \operatorname{Ext}^{<0}(E_v,E_v)=0.
\]
Its whole Chern character remains Hodge along the signature-$(4,4)$ unitary domain. Nevertheless its first-order obstruction map on that domain has rank six. No infinitesimal gerbe twist removes those six normal obstruction directions. For every finite translation and degree-zero-twist action under which $E_v$ is equivariant, they form a six-dimensional subspace of the kernel of the invariant semiregularity map.
\end{proposition}

The representative in the proposition is the direct sum of the graph structure sheaves in the signed expansion of $v$, with a shift by one on the negative summands. The proposition concerns this representative and its translates, degree-zero twists, and finite equivariant inductions. It is not a nonexistence assertion for every perfect complex with a Weil component.

\begin{proof}
Choose coordinates $x_1,\ldots,x_4$ and $y_1,\ldots,y_4$ on the two tangent spaces. The polarized $K$-linear deformation tangent space consists of
\begin{equation}\label{eq:unitary-tangent}
 \kappa_B=\begin{pmatrix}0&B\\B^{\mathsf t}&0\end{pmatrix},
 \qquad B\in M_4(\C).
\end{equation}
Indeed commuting with the complex multiplication action kills the two diagonal Beltrami blocks; compatibility with the polarization makes the full matrix symmetric. The resulting space has dimension $16$.

For a graph with tangent equation $y=Tx$, the normal quotient is $(v_x,v_y)\mapsto v_y-Tv_x$. Restricting \eqref{eq:unitary-tangent} to its antiholomorphic tangent and taking this quotient gives the normal obstruction
\[
 \nu_T(B)=B^{\mathsf t}-TB\bar T
 \in H^1(\Gamma_T,N_{\Gamma_T/A_0}).
\]
In particular, for the diagonal and for every scalar graph of norm-one slope,
\begin{equation}\label{eq:graph-normal-obstruction}
 \nu_{\lambda I}(B)=B^{\mathsf t}-B,
 \qquad \lambda\bar\lambda=1.
\end{equation}
Its kernel is the ten-dimensional space of symmetric matrices and its rank is six. The matrix $B=E_{12}$ gives the nonzero example $E_{21}-E_{12}$.

We check that this is an obstruction for the structure sheaf itself. The automorphism $(x,y)\mapsto(x,y-x)$ sends the diagonal to $X\times\{0\}$. K\"unneth and the Koszul resolution of the point of $Y$ give
\begin{equation}\label{eq:diagonal-ext}
 \operatorname{Ext}^2(\mathcal O_\Delta,\mathcal O_\Delta)
 =H^2(X,\mathcal O_X)
 \oplus\bigl(H^1(X,\mathcal O_X)\otimes T_{0,Y}\bigr)
 \oplus\bigwedge^2T_{0,Y}.
\end{equation}
The dimensions are $6,16,6$. Huybrechts--Thomas \cite[Corollary~3.4, manuscript p.~14]{HT} identify the obstruction of a perfect complex as the product of its Atiyah class with the Kodaira--Spencer class. In normal coordinates $u=y-x$, the Koszul Atiyah class has normal symbol $\sum_jdu_j\otimes e_j$. Its contraction is
\[
 \operatorname{ob}_{\kappa_B}(\mathcal O_\Delta)
 =\sum_{j,p}(B^{\mathsf t}-B)_{jp}\,
 d\bar x_p\otimes e_j,
\]
in the middle summand of \eqref{eq:diagonal-ext}. These constant harmonic representatives are independent. An infinitesimal gerbe class $\beta\in H^2(A_0,\mathcal O_{A_0})$ acts as $\beta|_\Delta\operatorname{id}$ in the first summand. It therefore cannot cancel a nonzero normal component.

For a rational norm-one slope, Hilbert~90 writes $\lambda=\bar c/c$ after clearing denominators. The $K$-endomorphism $(x,y)\mapsto(cx,\bar c y)$ takes the diagonal isogenously to that graph and extends over the unitary family. Pullback to the corresponding finite \'etale cover reduces the obstruction to the calculation above. Translations and tensoring by degree-zero line bundles preserve it.

The diagonal detects the Weil plane: a generator of one complex Weil line is
\[
 dx_1\wedge\cdots\wedge dx_4
 \wedge d\bar y_1\wedge\cdots\wedge d\bar y_4,
\]
whose restriction to $\Delta$ is a nonzero volume form. Its tangent and normal bundles are trivial, so Grothendieck--Riemann--Roch gives $\operatorname{ch}(\mathcal O_\Delta)=[\Delta]$, with no other components. On the five paired $K$-weights in degree eight, $S$ has eigenvalues
\[
 5^j\operatorname{Tr}_{K/\Q}(a^{8-2j}),\quad0\le j\le4,
 \qquad (-1054,-1170,-350,750,1250).
\]
The first is the Weil weight. Hence $P(S)$ kills the other weights and acts there by
$P(-1054)=-339429556224$.

Every isogeny pullback of $\mathcal O_\Delta$ is the structure sheaf of a reduced disjoint union of torsion translates of scalar norm-one graphs. Expand $v=\sum_Zm_Z[\mathcal O_Z]$, combining identical supports, and define
\[
 E_v=\bigoplus_{m_Z>0}\mathcal O_Z^{\oplus m_Z}
 \ \oplus\!
 \bigoplus_{m_Z<0}\mathcal O_Z[1]^{\oplus(-m_Z)}.
\]
This proves the Chern-character assertion. For distinct integral supports of the same dimension, $\operatorname{Hom}(\mathcal O_Z,\mathcal O_{Z'})=0$: a nonzero image inside the torsion-free rank-one sheaf $\mathcal O_{Z'}$ would have generic support $Z'$, forcing $Z'=Z$. These Hom groups are the only possible negative self-Ext groups introduced by the shifts zero and one. Thus $E_v$ is gluable, meaning that its negative self-Ext groups vanish.

Its Atiyah obstruction is block diagonal, with the nonzero normal classes \eqref{eq:graph-normal-obstruction} on every surviving summand. A shift changes the Chern-character trace but does not cancel different diagonal blocks in self-Ext. Since at least one coefficient $m_Z$ is nonzero, the obstruction map still has rank six. The preceding decomposition also shows that a gerbe class cannot cancel it.

Finally the Buchweitz--Flenner semiregularity map is the trace of the powers of the Atiyah class \cite[pp.~136--137]{BF}. With compatible sign conventions its components satisfy
\[
 \sigma_E^{2,q}(\operatorname{ob}_{\kappa}(E))
 =\kappa\mathbin{\lrcorner}\operatorname{ch}_{q+1}(E).
\]
This follows by expanding the exponential defining the Chern character and using cyclicity of trace. The only nonzero component of $\operatorname{ch}(E_v)$ is a flat Weil class, which remains Hodge along the whole unitary domain. Every displayed contraction therefore vanishes. The six-dimensional obstruction image is in the kernel of $\sigma_{E_v}^2$.

Translations preserve the constant tensors $\kappa_B$, and degree-zero line bundles have zero Atiyah class. Naturality makes these obstruction classes invariant under every compatible finite action of $G\subset A_0\times A_0^\vee$. A multiplicity representation fixes the identity endomorphism in each block, and a permutation of supports permutes the compatible obstruction blocks. Thus the same six-dimensional subspace lies in the kernel on $\operatorname{Ext}^2(E_v,E_v)^G$. Equivariant induction adds such blocks and does not remove them.
\end{proof}

In Perry's theorem the representative also fails the weak equivariant condition. Indeed, with $B_0=0$, all its other hypotheses are satisfied. Weak equivariant semiregularity would give a twisted deformation over an \'etale neighborhood in the unitary family. An \'etale base change preserves the tangent space, contradicting the nonzero obstruction even after an arbitrary infinitesimal gerbe twist. This last conclusion uses the semiregularity theorem in \cite{Perry}; the explicit obstruction calculation above uses the published Atiyah and semiregularity formulas.

Products of fourfold objects face a related issue. For $F\in D(A_1)$ and $H\in D(A_2)$, the Atiyah class of $F\boxtimes H$ is the sum of the two exterior Atiyah classes. A mixed Kodaira--Spencer direction $\eta_1\otimes v_2\in H^1(\mathcal O_{A_1})\otimes H^0(T_{A_2})$ therefore gives
\[
 (\eta_1\operatorname{id}_F)\boxtimes
 (v_2\mathbin{\lrcorner}\operatorname{At}(H))
 \in\operatorname{Ext}^1(F,F)\otimes\operatorname{Ext}^1(H,H).
\]
Semiregularity on the two factors controls the product locus and does not by itself remove this mixed obstruction. A construction for the full eightfold domain must address these extra directions.
